Protože $\sigma=n|q|\mu$, teplota ovlivňuje vodivost změnou koncentrace nosičů $n$ a jejich pohyblivosti $\mu$.
\begin{itemize}
\item \textbf{Kapalné izolanty:} rozhodující je změna pohyblivosti iontů. S teplotou klesá viskozita $\eta$, takže roste $\mu$ i vodivost; při neměnné koncentraci platí přibližně $\sigma\eta=\mathrm{konst.}$
\item \textbf{Pevné izolanty:} rozhodující je teplotně podmíněné uvolňování nosičů a vznik poruch mřížky. Vyšší teplota usnadňuje překonání energetických bariér. Ve skle se navíc mění pohyblivost iontů podle viskozity skelné sítě.
\end{itemize}
V přednášce je závislost vyjádřena $\sigma=Ae^{-b/T}$, kde $A,b$ jsou materiálové konstanty: s teplotou konduktivita roste.
